\chapter{Závěr}
\label{chap:conclusion}

V této práci byl úspěšně demonstrován a statisticky vyhodnocen vliv tokenizace, velikosti dat a kapacity modelu na jazykové modelování v nízkozdrojovém režimu minimalistického jazyka toki pona.

Výsledky ukazují, že:
\begin{enumerate}
    \item Celoslovní tokenizace (Word) poskytuje v malých datech nejvyšší datovou efektivitu a dosahuje až 100\,\% přesnosti v gramatických testech, neboť eliminuje zátěž spojenou se syntézou morfologických jednotek z jednotlivých znaků.
    \item Znaková tokenizace (Char) vykazuje nejstrmější škálovací exponent ($\beta = 0{,}238$), což znamená, že s rostoucím množstvím dat nejrychleji snižuje ztrátu.
    \item Chování modelů spolehlivě odpovídá mocninným škálovacím zákonům (Kaplan / Chinchilla) i na škále stovek tisíc parametrů.
\end{enumerate}
V budoucím výzkumu je vhodné rozšířit experimentální prostor na vícejazyčné nízkozdrojové korpusy a hlubší architektury.
